\section{Edge dropping}\label{sec:edge}
The classical edge rule of Henson, Lal and Pusey \cite{HLP,Evans,AWS} removes $a\to b$ under two hypotheses:
\begin{description}[leftmargin=1.2cm]
\item[(H1)] $a$ belongs to at least one source, and every source of $a$ also contains $b$.
\item[(H2)] $b$ sees every visible parent of $a$.
\end{description}
Classically, $b$ can compute $a$'s response from its settings, latents and private noise (a shared
latent can supply that noise). A quantum share generally cannot predict the measurement on another
share. Rules A and B give two ways for the target to recover the missing outcome.

\begin{figure}[H]
\centering
\begin{tikzpicture}[mdag]
  \begin{scope}
    \vn{y}{1.1,1.7}{$y$} \vn{a}{0,0.3}{$a$} \vn{b}{2.2,0.3}{$b$} \qs[csrc]{q}{1.1,-0.75}
    \draw[cedge] (q) -- node[lab,below left=-1pt,text=ccol] {$\lambda,\nu$} (a);
    \draw[cedge] (q) -- node[lab,below right=-1pt,text=ccol] {$\lambda,\nu$} (b);
    \de{y}{a} \de{y}{b} \de[dropedge]{a}{b}
    \node[tag,align=center] at (1.1,-1.75) {classical: the latent carries the noise $\nu$ of $a$;\\
      $b$ computes $x_a=f(y,\lambda,\nu)$};
  \end{scope}
  \begin{scope}[shift={(7.6,0)}]
    \vn{y}{1.1,1.7}{$y$} \vn{a}{0,0.3}{$a$} \vn{b}{2.2,0.3}{$b$} \qs{q}{1.1,-0.75}
    \feed{q}{a,b} \de{y}{a} \de{y}{b} \de{a}{b}
    \node[lab,text=dropcol] at (1.1,0.55) {?};
    \node[tag,align=center] at (1.1,-1.75) {quantum: $x_a$ does not exist before $a$ measures;\\
      the share of $b$ is another system, not a copy};
  \end{scope}
\end{tikzpicture}
\caption{Why the classical argument for dropping $a\to b$ has no quantum counterpart. On the left the source is
classical (blue) and hands $a$ and $b$ the same data.}
\label{fig:edge-idea}
\end{figure}

\subsection{Rule A: the origin and its settings}

\begin{rulebox}{Edge dropping, Rule A}
\begin{theorem}[Rule A: the origin is premeasurable]\label{thm:ruleA}
Let the edge $a\to b$ satisfy (H1) and (H2), and let $a$ be premeasurable. Then
\[\Q(\G)=\Q(\G-(a\to b))\qquad\text{and}\qquad\C(\G)=\C(\G-(a\to b)).\]
\end{theorem}
\end{rulebox}

\begin{proof}
By \cref{lem:premeasure}, $b$ can read $x_a$ from the shared source and select its old measurement.
Adding the edge back is structural dominance. Both constructions preserve diagonal models.
\end{proof}

\begin{figure}[H]
\centering
\begin{tikzpicture}[mdag]
  \begin{scope}
    \vn[pre]{a}{0,0}{$a$} \vn{b}{2.4,0}{$b$} \vn{c}{1.2,1.9}{$c$} \vn{t}{2.9,1.9}{$t$}
    \qs{q}{1.2,0.85}
    \feed{q}{a,b,c} \de[dropedge]{a}{b} \de{t}{c}
  \end{scope}
  \node[rel] at (3.9,0.8) {\equivsign};
  \begin{scope}[shift={(5.0,0)}]
    \vn[pre]{a}{0,0}{$a$} \vn{b}{2.4,0}{$b$} \vn{c}{1.2,1.9}{$c$} \vn{t}{2.9,1.9}{$t$}
    \qs{q}{1.2,0.85}
    \feed{q}{a,b,c} \de{t}{c}
  \end{scope}
  \node[tag,anchor=west] at (8.7,1.1) {because the source can measure};
  \node[tag,anchor=west] at (8.7,0.75) {$a$'s share and tell $b$ the result};
  \begin{scope}[shift={(8.9,-0.95)}]
    \vn[pre]{a}{0,0}{$a$} \vn{b}{2.4,0}{$b$} \qs{q}{1.2,0.85}
    \draw[cedge] (q) -- node[lab,above left,text=ccol] {$x_a$} (a);
    \draw[cedge] (q) to[bend left=20] node[lab,above right,text=ccol,pos=0.5] {$x_a$} (b);
    \draw[qedge] (q) to[bend right=20] (b);
  \end{scope}
\end{tikzpicture}
\caption{Rule A. The yellow origin has no setting and a single source, which also contains the target. The
bystander $c$ is arbitrary; here it has a setting.}
\label{fig:ruleA}
\end{figure}

\begin{proposition}[Rule A with settings that every member sees]\label{prop:ruleA-settings}
Let the edge $a\to b$ satisfy (H1), let $a$ be a member of exactly one source $F$, and let every setting of $a$ be
a setting of every other member of $F$. Then $\Q(\G)=\Q(\G-(a\to b))$ and $\C(\G)=\C(\G-(a\to b))$.
\end{proposition}

\begin{proof}
For every setting $y$ the sole source $F$ prepares an independent entry
\[
 \omega_y=\sum_x\proj{x}_{R_a}\otimes\proj{x}_{R_b}\otimes
          \Tr_{F,a}\!\left[(M_a^{x\mid y}\otimes\id)\rho_F\right].
\]
All members know $y$. In that entry $a$ outputs its register and $b$ uses its copy to select its
old measurement. Every other member measures its share as before. The selected term in \cref{eq:trace}
is the original one and the unused entries have trace one. Diagonal models stay diagonal.
\end{proof}

\begin{validbox}[label=ex:target-two]{Rule A; the target may have further sources}
\begin{minipage}{0.55\linewidth}
The origin $0$ is premeasurable with the source $\{0,1\}$. The target $1$ is also a member of $\{1,2\}$ and may
measure its two shares jointly; Rule A does not mind:\\[3pt]
\str{0→1; Q\{0,1\}, Q\{1,2\}} \ $\equiv$ \ \str{Q\{0,1\}, Q\{1,2\}}.\\[3pt]
The reverse move adds an edge from a premeasurable setting in \cref{fig:bell-class}.
\end{minipage}\hfill
\begin{minipage}{0.42\linewidth}\centering
\begin{tikzpicture}[mdag,scale=0.85]
  \newcommand{\chainnodes}{\vn[pre]{a}{0,0}{$0$} \vn{b}{1.5,0}{$1$} \vn{c}{3.0,0}{$2$}
    \qs{q}{0.75,0.9} \qs{r}{2.25,0.9} \feed{q}{a,b} \feed{r}{b,c}}
  \chainnodes \de[dropedge]{a}{b}
  \node[rel] at (1.5,-0.8) {\equivsign};
  \begin{scope}[shift={(0,-2.5)}] \chainnodes \end{scope}
\end{tikzpicture}
\end{minipage}
\end{validbox}

\subsection{Rule B: the bystanders}
The bystanders are the members other than $a,b$ of the sources containing $a$. The origin can now
have several sources and its measurements can depend on settings.

\begin{rulebox}{Edge dropping, Rule B}
\begin{theorem}[Rule B: every bystander is premeasurable]\label{thm:ruleB}
Let the edge $a\to b$ satisfy (H1) and (H2), and let every bystander be premeasurable. Then
\[\Q(\G)=\Q(\G-(a\to b))\qquad\text{and}\qquad\C(\G)=\C(\G-(a\to b)).\]
\end{theorem}
\end{rulebox}

\begin{figure}[H]
\centering
\begin{tikzpicture}[mdag]
  \newcommand{\ruleBnodes}{\vn{y}{1.3,4.2}{$y$} \vn{a}{0,1.0}{$a$} \vn{b}{2.6,1.0}{$b$}
    \vn[pre]{u}{1.3,-1.0}{$u$} \vn[pre]{up}{1.3,3.0}{$u'$}
    \qs{f1}{1.3,-0.05} \qs{f2}{1.3,2.05}
    \feed{f1}{a,b,u} \feed{f2}{a,b,up}
    \draw[edge] (y) to[bend right=28] (a); \draw[edge] (y) to[bend left=28] (b);}
  \begin{scope}
    \ruleBnodes \de[dropedge]{a}{b}
    \node[lab,text=qcol] at (1.78,-0.4) {$F_1$}; \node[lab,text=qcol] at (1.78,2.38) {$F_2$};
  \end{scope}
  \node[rel] at (3.7,1.0) {\equivsign};
  \begin{scope}[shift={(4.8,0)}]
    \ruleBnodes
  \end{scope}
  \begin{scope}[shift={(9.1,0.7)}]
    \tikzset{cell/.style={rounded corners=1.5pt, minimum width=10mm, minimum height=4.4mm, font=\scriptsize}}
    \node[tag,anchor=west] at (-0.3,3.0) {the catalogue that $F_1$ emits:};
    \node[lab,anchor=south] at (1.75,2.2) {outcome}; \node[lab,anchor=south] at (3.2,2.2) {state};
    \foreach \j/\lbl/\w in {0/{$(y,\ell)=(0,0)$}/0.5pt, 1/{$(y,\ell)=(0,1)$}/0.5pt, 2/{$(y,\ell)=(1,0)$}/1.1pt, 3/{$(y,\ell)=(1,1)$}/0.5pt} {
      \node[lab,anchor=east] at (1.05,{1.9-0.58*\j}) {\lbl};
      \node[cell, draw=ccol, fill=ccol!10, line width=\w] at (1.75,{1.9-0.58*\j}) {$\xi_{y,\ell}$};
      \node[cell, draw=qcol, fill=qcol!12, line width=\w] at (3.2,{1.9-0.58*\j}) {$\sigma$};
    }
    \node[lab,anchor=north] at (1.75,-0.22) {to $a$ and $b$}; \node[lab,anchor=north] at (3.2,-0.22) {to $b$};
    \node[tag,anchor=west] at (-0.3,-1.1) {$a$ and $b$ both know $y$ and $\ell$,};
    \node[tag,anchor=west] at (-0.3,-1.45) {so both use the same entry};
  \end{scope}
\end{tikzpicture}
\caption{Rule B. The origin $a$ has a setting $y$, which $b$ also sees, and two sources, both of which contain
$b$. The bystanders $u$ and $u'$ are premeasurable; $\ell$ stands for their outcomes.}
\label{fig:ruleB}
\end{figure}

\begin{proof}
Write $F_i$ for the sources of $a$ and $U_i=F_i\setminus\{a,b\}$. Pre-measure each bystander.
Source $F_i$ draws its bystander outcomes $\ell_i$ and shares them with its members. Before measuring
$a$, these draws are independent across sources, with joint weight $p(\ell)=\prod_i p_i(\ell_i)$.
For fixed $\ell$ and origin settings $y$, measuring $a$ on all its shares gives a distribution
$r(x\mid y,\ell)$ and a conditional state $\sigma_{x,y,\ell}$ on $b$'s shares of the $F_i$.
These are the only systems left in those sources. In the original model their unnormalised state is
\[
 \tau(x;y,\ell)=\Tr_{a,U}\!\left[
  \left(M_a^{x\mid y}\otimes\bigotimes_{u\in U}M_u^{\ell_u}\otimes\id_b\right)
  \bigotimes_i\rho_{F_i}\right]
 =p(\ell)r(x\mid y,\ell)\sigma_{x,y,\ell},
\]
where $U=\bigcup_i U_i$ and the trace is over the shares held by $a$ and $U$.
For a zero-weight entry one may choose any fixed state.

Have one shared source prepare a catalogue indexed by every $(y,\ell)$, independently of the draws.
In each entry it samples $x$ from $r$, gives $a,b$ a copy of $x$, and gives $b$ the state
$\sigma_{x,y,\ell}$. Both nodes see all the $\ell_i$ by (H1), and both see $y$ by (H2).
Node $a$ outputs the selected $x$; $b$ uses it to measure the selected state jointly with its other shares.
For each visible assignment the weight $p(\ell)r(x\mid y,\ell)$ times the conditional state recovers
the original unnormalised state at $b$. Thus \cref{eq:trace} agrees, including on zero-weight assignments.
The same construction is classical for diagonal inputs, and the reverse inclusion is structural dominance.
\end{proof}

\begin{validbox}[label=ex:ruleB-bystander]{Rule B; a premeasurable bystander}
\begin{minipage}{0.55\linewidth}
The origin $1$ has the setting $0$, seen by the target $2$; the source $\{1,2,3\}$ has the bystander $3$, which has
no setting and no other source:\\[3pt]
\str{0→1, 0→2, 1→2; Q\{1,2,3\}} \ $\equiv$ \ \str{0→1, 0→2; Q\{1,2,3\}}.\\[3pt]
Giving $3$ the setting $0$ makes every other member of the source see the origin's settings. The wider Rule A
then drops the edge (\cref{prop:ruleA-settings}), while the setting at $3$ rules out Rule B.
\end{minipage}\hfill
\begin{minipage}{0.42\linewidth}\centering
\begin{tikzpicture}[mdag,scale=0.85]
  \vn{y}{0.9,1.5}{$0$} \vn{a}{0,0}{$1$} \vn{b}{1.8,0}{$2$} \vn[pre]{o}{0.9,-1.6}{$3$} \qs{q}{0.9,-0.7}
  \feed{q}{a,b,o} \de{y}{a} \de{y}{b} \de[dropedge]{a}{b}
  \node[rel] at (2.6,0.0) {\equivsign};
  \begin{scope}[shift={(3.4,0)}]
    \vn{y}{0.9,1.5}{$0$} \vn{a}{0,0}{$1$} \vn{b}{1.8,0}{$2$} \vn[pre]{o}{0.9,-1.6}{$3$} \qs{q}{0.9,-0.7}
    \feed{q}{a,b,o} \de{y}{a} \de{y}{b}
  \end{scope}
\end{tikzpicture}
\end{minipage}
\end{validbox}

\paragraph{Using the rules together.}\label{sec:compose}
A premeasurable $p$ can send its outcome along a new edge to each other member of its sole source $F$
by Rule A. Face splitting with $C=\{p\}$ then replaces $F$ with $F\setminus\{p\}$.
The Bell examples use this exchange in both directions.

For example, add the edges from $p$ to the other members of $\{p,c,d\}$. Now $c,d$ both know $x_p$
and the source reduces to $\{c,d\}$. If a remaining source already contains $c,d$, it can carry
their catalogue. The visible arrows from $p$ still carry its outcome.

\begin{lemma}\label{lem:simultaneous}
Let $\Delta$ be the set of edges of $\G$ that satisfy Rule A with shared settings (\cref{prop:ruleA-settings})
or Rule B (\cref{thm:ruleB}) in $\G$. Then
$\G-\Delta\qeq\G$ and $\G-\Delta\ceq\G$.
\end{lemma}
\begin{proof}
Fix a topological order of $\G$ and remove the edges $(a,b)\in\Delta$ one at a time, those with a later origin $a$
first. Sources stay fixed, so (H1) persists. Removing edges also preserves the premeasurability of every
bystander in Rule B. Let $y$ be any remaining setting of $a$ when $a\to b$ is removed. For (H2) the original
graph contains $y\to b$. For Rule A it also contains $y\to v$ for every other member $v$ of the sole source of
$a$. All these edges are still present because their origin $y$ precedes $a$ in the topological order.
\end{proof}

Widening Rule A to shared settings licenses $35$ additional four-node edge deletions within known
classes. On five nodes it joins $542$ classes, reducing the count from $596{,}413$ to $595{,}871$.
